\documentclass[a4paper,12pt,twocolumn]{article}

\usepackage[landscape,top=1cm,bottom=1.4cm,left=1cm,right=1cm,columnsep=1cm]{geometry}
\usepackage[fleqn]{amsmath}
\usepackage{amssymb}
\usepackage{fontspec}
\usepackage{enumitem}
\usepackage{xeCJK}
\usepackage[hidelinks]{hyperref}

\setlength{\mathindent}{1.5em}

% Set fonts
\setmainfont{Times New Roman}
\setsansfont{Arial}
\setmonofont{JetBrains Mono}

% Chinese font support
\setCJKmainfont{Iansui}[
  BoldFont=Iansui,
  AutoFakeBold=3.17
]

\pagestyle{plain}
\setlength{\columnseprule}{0.2pt}

\newcommand{\examdate}{2026/03/26}
\newlength{\problemnumberwidth}
\settowidth{\problemnumberwidth}{\textbf{94.}\ }
\newcommand{\problemtitle}[2]{%
  \hypertarget{#1}{}%
  \par\vspace{0.4em}%
  \noindent{\large\bfseries #2}\par
  \vspace{0.3em}%
  \hrule
  \vspace{0.9em}%
}
\newcommand{\worksheettoc}{%
  \noindent{\Large\bfseries Contents}\par
  \vspace{0.3em}%
  \hrule
  \vspace{0.9em}%
  \begin{tabular}{@{}p{0.26\columnwidth}p{0.66\columnwidth}@{}}
    \hyperlink{sec:10-1}{\textbf{10-1}} & \hyperlink{sec:10-1}{16, 19, 22, 57} \\
    \hyperlink{sec:10-3}{\textbf{10-3}} & \hyperlink{sec:10-3}{62} \\
    \hyperlink{sec:11-3}{\textbf{11-3}} & \hyperlink{sec:11-3}{34} \\
    \hyperlink{sec:11-10}{\textbf{11-10}} & \hyperlink{sec:11-10}{1, 2, 93, 94} \\
  \end{tabular}\par
}
\newcommand{\problem}[2]{%
  \par\noindent
  \hangindent=\problemnumberwidth
  \hangafter=1
  \makebox[\problemnumberwidth][l]{\textbf{#1.}}#2\par
  \vspace{1.4em}%
}
\newlist{subparts}{enumerate}{1}
\setlist[subparts]{%
  label=(\alph*),
  labelindent=\problemnumberwidth,
  leftmargin=*,
  labelsep=0.5em,
  align=left,
  itemsep=0.2em,
  topsep=0.3em,
  parsep=0pt,
  partopsep=0pt
}

\begin{document}

\twocolumn[
{\centering
\Large\bfseries Calculus Problems\par
\vspace{0.5em}
\noindent\hfill\normalsize\normalfont \examdate\par
\vspace{0.8em}
}
]

\worksheettoc
\vfill\null
\newpage

\problemtitle{sec:10-1}{10-1}

\problem{16}{%
\begin{subparts}
\item Eliminate the parameter to find a Cartesian equation of the curve.
\item Sketch the curve and indicate with an arrow the direction in which the curve is traced as the parameter increases.
\end{subparts}
\[
x=\csc t,\ y=\cot t,\ 0<t<\pi
\]}

\problem{19}{%
\begin{subparts}
\item Eliminate the parameter to find a Cartesian equation of the curve.
\item Sketch the curve and indicate with an arrow the direction in which the curve is traced as the parameter increases.
\end{subparts}
\[
x=\ln t,\ y=\sqrt{t},\ t\geq 1
\]}

\problem{22}{%
\begin{subparts}
\item Eliminate the parameter to find a Cartesian equation of the curve.
\item Sketch the curve and indicate with an arrow the direction in which the curve is traced as the parameter increases.
\end{subparts}
\[
x=\sinh t,\ y=\cosh t
\]}

\problem{57}{Find the point at which the parametric curve intersects itself and the corresponding values of \(t\).
\begin{subparts}
\item \(x=1-t^2,\ y=t-t^3\)
\item \(x=2t-t^3,\ y=t-t^2\)
\end{subparts}}

\problemtitle{sec:10-3}{10-3}

\problem{62}{Graph the polar curve. Choose a parameter interval that produces the entire curve.
\[
r=\lvert\tan\theta\rvert^{\lvert\cot\theta\rvert}
\]
\noindent(valentine curve)}

\problemtitle{sec:11-3}{11-3}

\problem{34}{Find the values of \(p\) for which the series is convergent.
\[
\displaystyle \sum_{n=1}^{\infty}\frac{\ln n}{n^p}
\]}

\problemtitle{sec:11-10}{11-10}

\problem{1}{Use the definition of a Taylor series to find the first four nonzero terms of the series for \(f(x)\) centered at the given value of \(a\). (6 point)
\[
f(x)=\sqrt[3]{x},\ a=8
\]}

\problem{2}{Evaluate the indefinite integral as an infinite series. (7 point)
\[
\displaystyle \int \frac{\cos x-1}{x}\,dx
\]}

\problem{93}{Use the Maclaurin series for \(f(x)=x\sin(x^2)\) to find \(f^{(203)}(0)\).
\medskip
\noindent\textit{Note:} This problem uses the Maclaurin series, so the expansion is centered at \(0\). If the center were changed to some nonzero value \(a\), is there still a way to solve it by using a Taylor series centered at \(a\)?}

\problem{94}{If \(f(x)=(1+x^3)^{30}\), what is \(f^{(58)}(0)\)?}

\end{document}
